
Las pruebas de usabilidad consisten en analizar la conducta del usuario frente a un artefacto para detectar posibles problemas y proponer soluciones. La prueba se puede hacer cuando el artefacto está en fase de diseño, en cuyo caso se realiza la prueba sobre un prototipo, o cuando está ya diseñado y siendo utilizado por el usuario final \cite{PUsabilidad}.

Por otro lado, las pruebas cognitivas identifican y miden los niveles específicos del rendimiento de la persona en algunas áreas esenciales como: lectura, matemáticas, lenguaje escrito y comunicación. Si bien estas pruebas pretenden determinar la posible existencia de un problema en cualquiera de estas áreas, es responsabilidad directa del evaluador analizar los resultados en función de las habilidades que el público objetivo posee. \cite{PCognitivas}

De este modo las pruebas de usabilidad y las pruebas cognitivas planteadas por los autores se representan mediante los test's, donde se buscan identificar las perspectivas singulares de cada uno de los estudiantes al implementar el sistema hipermedia con características adaptativas, de este modo evaluando la usabilidad y los procesos cognitivos que estuvieron implicados en el desarrollo de alguna tarea. A continuación, se visualizará las etapas que ayudaron a la identificación de los problemas del sistema; como sus resultados.

\section{Elaboración del problema.}

Se ha formulado el problema con una pregunta concisa y crucial dentro de la implementación del sistema hipermedia con características adaptativas: ¿Qué características de usabilidad y estrategias cognitivas – emocionales están presentes en el sistema hipermedia con características adaptativas para los estudiantes de la Institución Educativa Juan María Céspedes?

\section{Definición de la población objetivo.}

La población objetivo será los estudiantes de grado 10 de la Institución Educativa Juan María Céspedes del municipio de Tuluá - Valle

\section{Diseño de la muestra.}

Los test’s se han diseñado para un grupo del grado 10$^{\circ}$ de la Institución Educativa Juan María Céspedes, no obstante, se ha decidido abordar una muestra de 20 estudiantes de 157 estudiantes registrados en la Institución, puesto que, la Institución no cuenta con los equipos de cómputo necesarios para abordar una mayor muestra de estudiantes.

No obstante, estas limitaciones no son un problema para las pruebas de usabilidad puesto que según el libro ``No me hagas pensar: Una aproximación a la usabilidad en la web'' \cite{NomeHagasPensar}, el autor declara que ``el secreto mejor guardado de la prueba de usabilidad es el grado al que no importa a quien pruebe ni cuántos usuarios prueban la aplicación'', de lo anterior, se puede verificar que la muestra es apropiada para realizar las pruebas de usabilidad para identificar las dificultades existentes en el aplicación.

\section{Test de usabilidad.}

La usabilidad es una disciplina que está presente en los sistemas computacionales, esta disciplina surge debido a la implementación de aplicaciones de uso comercial en internet; el área que busca apoyar la usabilidad es el desarrollo y diseño de interfaces computacionales que faciliten de una u otra manera las tareas que las personas llevan frente a una pantalla.

Para que el sistema computacional cumpla con las características de la usabilidad se establecen estándares generales a través de los cuales es posible conseguir que el usuario tenga una experiencia positiva frente a los sistemas computacionales. No obstante, es importante considerar que, la usabilidad siempre busca el mismo objetivo, puesto que su propósito fundamental es lograr que el usuario pueda interactuar con el contenido y las funciones del sistema computacional de manera intuitiva, simple y directa, no obstante, los elementos que se miden dentro del sistema computacional pueden variar dependiendo de las características singulares que tengan las interfaces, por otro lado, hay normas generales, como las que plantean Jakob  Nielsen\cite{HJakob}, de las cuales se pueden desprender normas singulares para un determinado tipo de sistema.

Se planteó un test de usabilidad el cual permitió determinar cuáles son las tareas más difíciles de completar por parte de los estudiantes, así como los elementos de la aplicación  que sean menos comprensibles. Esa información se evaluó y priorizó con el objetivo de mejorar el sistema hipermedia con características adaptativos en el ámbito de que la aplicación pueda ser usable para estudiantes.

El test se estructuró para comprender 3 secciones con el fin de conocer el perfil de los estudiantes que realizan dicho test, se ha estructurado de la siguiente manera: Pre - Test, Test de Usabilidad y Post- Test; No obstante, se tuvo en cuenta la opinión de los docentes de la Universidad del Valle y del coordinador del trabajo de grado en la estructuración de los test.

\textbf{Pre - Test:} Esta sección se realiza con el fin de identificar las características de los estudiantes de manera individual tales como, las horas que pasan navegando, su edad, y los sitios web favoritos, con el fin de establecer el perfil de la muestra, esta sección cuenta con 8 preguntas.

\textbf{Test de Usabilidad:} Esta sección se realiza con el propósito de identificar si los componentes ubicados de manera estratégica fueron los apropiados para el estudiante, si el estudiante logro identificar los elementos que conforman el sistema de manera individual, y si su interacción con ellos fue simple e intuitiva, así, se evitará la frustración por el hecho de que algo no está bien explicado o sea poco intuitivo. Las preguntas fueron estructuradas para medir diferentes ámbitos, cada uno de los cuales se analizan por separado.

\begin{itemize}

\item \textbf{Contenido}

Las preguntas de esta sección, se realiza luego de permitir al usuario navegar en la aplicación, con el fin de que se forme una opinión acerca de lo que está viendo y la forma de navegar por los contenidos (Temas y subtemas).

El objetivo de esta sección es determinar la calidad que proporcionan los contenidos y si la forma de presentarlos le permite al estudiante hacerse una idea concreta de la información que se le está entregando través de la aplicación. 

\item \textbf{Navegación} 

Las preguntas de esta sección permiten establecer si la forma de organizar la información dentro de la aplicación es adecuada de acuerdo a la experiencia, conocimientos y expectativas que tenga el estudiante que visita la aplicación

\item \textbf{Gráficos Web}

Las preguntas de esta sección buscan establecer si al estudiante le ayuda la información gráfica que se le ofrece a través de los temas y subtemas de la aplicación.

\item \textbf{Feedback}

Las preguntas de esta sección buscan establecer si el estudiante encuentra una forma de enviar formularios y que la aplicación responda positiva o negativamente dependiendo la acción que realiza el estudiante.

\item \textbf{Utilidad}

Las preguntas de esta sección son las finales de la prueba y tienen el objetivo de establecer una especie de resumen general de la experiencia de navegar en la aplicación.
\end{itemize}
\textbf{Post – Test:} Esta sección se realiza con el fin de determinar la percepción del estudiante frente al sistema en aspectos de contenido y visualización de imágenes después de que el estudiante interactuó con el sistema hipermedia con características adaptativas, esta sección cuenta con 5 preguntas.


\section{Test cognitivo.}

La psicología cognitiva también hace parte de los test’s que se implementaron en busca de mejorar aspectos funcionales y generales del sistema hipermedia con características adaptativas; puesto que, el estudiante asimila la información previamente suministrada y realiza los procesos mentales para el aprendizaje del contenido, se abordaron elementos cognitivos como el conocimiento, la reorganización y manifestación de datos, la percepción de los elementos del sistema de manera individual. Por lo anterior el test, se ha dividido en 6 secciones, estas secciones hacen parte de la taxonomía de Benjamin Bloom \cite{Benjamin}, a continuación, se visualizarán las características básicas de cada una de ellas.

\begin{itemize}
    \item \textbf{Recordar información previamente aprendida.}
    
    Son las características cognitivas de los estudiantes que permiten reconocer, reproducir, nombrar y describir cualquier información anteriormente adquirida, para identificar estas características se ha diseñado las preguntas N$^{\circ}$ 1 y 3.
    
    \item \textbf{Comprender.}
    
    Son las características cognitivas donde los estudiantes implementan para entender, apropiarse de aquellos temas que se han visto durante periodos de tiempo o de enseñanza, para identificar estas características se ha diseñado las preguntas N$^{\circ}$ 4, al 18.
    
    \item \textbf{Aplicar.}
    
    Son las características cognitivas donde el estudiante emplea o hace uso de un conocimiento en algunas situaciones o actividades que se presenten, para identificar estas características se ha diseñado las preguntas N$^{\circ}$ 19 y 20.
    
    \item \textbf{Evaluar.}
    
    Son las características cognitivas de los estudiantes para estimar o apreciar el valor de una situación singular,donde se deja alguna experiencia significativa, para identificar estas características se ha diseñado las pregunta N$^{\circ}$  21.
    
    \item \textbf{Analizar.}
    
    Son las características cognitivas de los estudiantes que implementan para distinguir las partes de un todo y relacionarlas con su estructura global.para identificar estas características se ha diseñado las preguntas N$^{\circ}$ 2 y 22.
        \end{itemize}

La taxonomía de Bloom guió a los docentes durante más de 40 años, y se considera uno de los textos educativos más significativos del siglo XX, antes de construir el esquema clasificador se debe tener claro qué es lo que se propone clasificar \cite{Benjamin}, partiendo de este hecho la taxonomía permitirá en el test cognitivo una clasificación de los procesos mentales y educativos de los estudiantes de la Institución Educativa Juan María Céspedes generados al implementar la aplicación con el propósito de poder identificar si el contenido es el apropiado y si el estudiante asimila los elementos presentes de forma intuitiva de acuerdo a sus procesos cognitivos.

\section{Test motivacional.}

Uno de los factores fundamentales que se debe tener en cuenta en el proceso de aprendizaje, es la motivación \cite{TesisBarrientos}. Siendo este factor, el que induce al deseo de aprender, y el esfuerzo por alcanzar metas, por lo tanto, no se puede pensar, que exista un aprendizaje exitoso si el estudiante no está motivado.

Si el estudiante no siente necesidad de aprender lo que el sistema hipermedia con características adaptativas le brinda, ni la necesidad de aplicar lo aprendido o de esforzarse y trabajar hasta sentirse satisfecho al resolver los cuestionarios del sistema, entonces lo que el sistema hipermedia brinda no tiene ningún valor, con lo anterior, se puede ver la importancia de valorar la motivación de los estudiantes del grado 10º de la Institución Educativa Juan María Céspedes en torno al sistema hipermedia evaluado para así determinar, si el factor de la motivación está afectando el aprendizaje del contenido presentado. Por lo anterior se ha planteado 7 preguntas.

\section{Preparación de la tarea de campo para la recolección de la información.}

Se realizó una solicitud a la secretaría de la Universidad del Valle, donde se expresa un permiso para llevar a cabo el test de usabilidad, test motivacional y el test cognitivo que es dirigida a la secretaría de la Institución Educativa Juan María Céspedes, la cual accede a brindar un espacio acordando con la coordinadora y los profesores de Física; no obstante, se procede a realizar la encuesta tomando una hora por cada grupo encuestado.


\section{Procesamiento y análisis de los datos.}

Para el procesamiento de datos se ha escogido la herramienta ofimática Excel, la cual se procedió a digitar las preguntas con sus debidas respuestas generadas por cada uno de los estudiantes y posterior a esto, se realizó su debida tabulación y gráfica estadística. De este hecho se analizó los test por cada tipología, describiendo su análisis de manera porcentual, textual, gráfica y sus debidas conclusiones por cada pregunta.

\section{Discusión de los resultados.}

Los resultados obtenidos se discuten en compañía del director de trabajo, para analizar los elementos que deben estar presente en el sistema hipermedia con características adaptativas con el objetivo de que cumplan con las necesidades del público objetivo, generando ideas constantes para el desarrollo.

Mediante las pruebas de usabilidad, cognitivas y motivacionales se realizó una aproximación del perfil de los estudiantes siendo así, el público objetivo que tiene el sistema hipermedia con características adaptativas, a continuación, se menciona los resultados de cada uno de los test’s.

Resultados del test de usabilidad: El test de usabilidad se dividió en tres secciones, las cuales fueron Pre – test, Test de usabilidad y el Post – Test.

\textbf{Pre – test:} En esta sección se abarcaron las preguntas correspondientes a la exploración del público objetivo, como se ha mencionado anteriormente; a continuación se visualizará las características singulares encontradas del público objetivo.
\begin{itemize}
    \item De acuerdo a los resultados de las preguntas N$^{\circ}$ 1, 2, 3, 4 y 7 se indican que el 85\% de los estudiantes encuestados navegan habitualmente en sitios web y pasan un promedio de 3,5 horas navegando, los sitios que visitan habitualmente para estudiar son las plataformas de consulta web (Wikipedia, WikiAprende, entre otros), el rango de edad son los 17 años, y tienen experiencia en la internet. 
    
    \item Los resultados de las preguntas N$^{\circ}$ 5, 6 y 8 indican una clara preferencia de los navegadores y la herramienta de búsqueda por parte de los estudiantes, donde los navegadores que más utilizan son: Google Chrome y Firefox Mozilla, con un 85\% y 10\% respectivamente, el 5\% indica otro navegador, pero no el nombre; la herramienta de búsqueda seleccionadas por los estudiantes en un 85\% es el buscador de Google. 
    
    \item Se puede concluir que la aplicación contempla las características del público objetivo puesto que, existe compatibilidad de los navegadores de Google Chrome y Mozilla y el contenido, elementos de navegación son acorde para el público objetivo.
    
\end{itemize}

\textbf{Test de usabilidad:} En esta sección se diseñó con el propósito de identificar si los componentes ubicados de manera estratégica fueron los apropiados para el público objetivo, como se ha mencionado anteriormente; a continuación se visualizará las características singulares encontradas del público objetivo.
\begin{itemize}

\item De acuerdo a los resultados de las preguntas N$^{\circ}$ 1, 2, 5 y 19 el 70\% de los estudiantes consideran que los textos usados en el contenido de los enlaces son suficientemente descriptivos y que no existen demasiados textos en estos, a su vez, el 85\% cree que el contenido sirven de apoyo para sus procesos de aprendizaje; por último el 80\% de los estudiantes consideran que el contenido fue planteado de forma secuencial, lo que resalta que hubo una buena secuenciación de los temas de la cinemática por parte de los autores. 

\item Los resultados de las preguntas N$^{\circ}$ 3, 4, 13, 14, 15 y 16 el 80\% de los estudiantes consideran que las imágenes implementadas en la aplicación fueron suficientemente descriptivas, el 90\% consideran que las imágenes presentadas en el contenido se relacionan con el texto, por otra parte, el 85\% de los estudiantes indican que la apariencia de la aplicación es la apropiada, es equilibrado en cuestiones de imágenes y  las imágenes fueron estructuradas de manera asertiva para servir de apoyo a los procesos de aprendizaje, debido a que, asemejan la pregunta con una imagen de manera fluida y directa.

\item Las preguntas N$^{\circ}$ 6, 7, 8, 9, 10, 11, 12 hacen referencia a la estructuración de los elementos de navegación dentro del sitio, considerando que el 80\% de los estudiantes, identificaron dichos elementos en tareas sencillas como llegar a un tema, en todo momento el sitio les indico en que subtema o tema se encontraban, y su implementación fue intuitiva; La pregunta 17 comprende las tareas de envío y notificación de formularios donde el 75\% de los estudiante afirman que la aplicación les confirmo cuando enviaron alguna información, mientras que el 10\% de los estudiantes afirman que el aplicación no les confirmo cuando enviaron información y el 15\% de los estudiantes no han indicado ninguna opción, por consiguiente se concluye que el aplicación interactúa con el estudiante no solo por el contenido, sino en la confirmación de los datos suministrados en algún formulario.

\item La pregunta N$^{\circ}$ 18 tiene como propósito determinar si el estudiante a primera vista, entiende el objetivo de la aplicación, lo que infiere que 18 estudiante de 20 estudiantes, han logrado identificar el objetivo de la aplicación, no obstante, se tendrá que replantear el inicio de la aplicación con texto alusivos con el fin de generar mayor claridad al respecto.

\end{itemize}

\textbf{Post – test:} En esta sección se abarcaron las preguntas correspondientes a la percepción del público objetivo de la aplicación, como se ha mencionado anteriormente; 
\begin{itemize}

\item  De acuerdo a los resultados de la pregunta N$^{\circ}$  1, se indica que solo el 5\% de los estudiantes encuestados consideran que la implementación de la aplicación ha sido difícil, mientras que el 95\% de los estudiantes consideran que la implementación de la aplicación está entre el rango de [Muy fácil (4) – Fácil (3) – Normal (7)], concluyendo que las características de usabilidad están presente en la aplicación. 

\item La pregunta N$^{\circ}$ 2, indica que el 75\% de los estudiantes consideran que encontrar una opción dentro de la aplicación esta entre el rango de [Muy fácil (4) – Fácil (4) – Normal (7)], mientras que el 25\% de los estudiantes no seleccionaron alguna opción. 

\item La pregunta N$^{\circ}$ 3, indica que el 65\% de los estudiantes consideran que comprender los cuestionarios de alguno de los temas ha estado presente entre el rango de [Muy fácil (1) – Fácil (2) – Normal (10)], mientras que el 10\% consideran que ha sido difícil comprender los cuestionarios y el 25\% no ha indicado opción alguna, concluyendo que la estructura y formulación de los cuestionarios es acorde a los procesos de aprendizaje que el estudiante ha desarrollado para la asignatura de Física. 

\item La pregunta N$^{\circ}$ 4 tiene como objetivo determinar si los estudiantes comprenden los ejercicios planteados en el contenido de la aplicación, lo que infiere que el 70\% de los estudiantes ha estado presente en el rango de [Muy fácil (1) – Fácil (8) – Normal (5)], mientras que el 5\% de los estudiantes consideran los ejercicios de los contenidos difíciles y el 25\% de los estudiantes no han seleccionado alguna opción, concluyendo que los ejercicios están bien estructurado de acuerdo a los procesos de aprendizaje del estudiante. 

\item La pregunta N$^{\circ}$ 5 tiene cómo objetivo principal determinar como ha sido la navegación de la aplicación, el 70\% de los estudiantes se encuentran en el rango de [Muy fácil (3) – Fácil (8) – Normal (3)], mientras que el 5\% ha indicado que ha sido difícil la navegación y el 25\% no han seleccionado respuesta alguna, concluyendo que la navegación ha sido planteada de acorde a las necesidades de los estudiantes de manera sencilla e intuitiva; para concluir también se puede evidenciar un porcentaje considerable de los estudiantes que no seleccionaron alguna respuesta, puesto que, indica que el público objetivo de alguna cierta manera no estuvo motivado al presentar el test o se le pasó por alto seleccionar las respuestas, de cualquier modo la aplicación cumple con las características de usabilidad de acuerdo a los resultados de esta sección.

\end{itemize}

\textbf{Test motivacional:} Se puede constatar que el perfil motivacional del estudiante posee características específicas como:
\begin{itemize}

\item Se puede evidenciar que en la pregunta N$^{\circ}$ 1 al estudiar, el estudiante lo hace para adquirir conocimiento o reforzar el conocimiento indicando que se tiene un objetivo principal lo cual lo impulsa a explorar y desarrollar los diferentes temas de la asignatura de Física para así alcanzar su objetivo.

\item Según las respuestas de la pregunta N$^{\circ}$ 2, el estudiante tiene muchas expectativas al empezar la clase de la asignatura de Física puesto que los rangos de las respuestas están en [Siempre (6) – Casi siempre (7) – Algunas veces (7)] concluyendo que el estudiante está atento frente a los temas de la asignatura

\item Se puede observar en la pregunta N$^{\circ}$ 3 que el estudiante es receptivo de manera positivo a los contenidos bien estructurado puesto que, la pregunta tiene como objetivo determinar si a mayor contenido de Física el estudiante considera que mejorará su desempeño, el rango evidenciado es [Muy de acuerdo (12) – De acuerdo (7) – Indiferente (1)]

\item Se puede evidenciar en la pregunta N$^{\circ}$ 4 que el estudiante en algunas ocasiones no opta por explicar la solución de los cuestionarios a los demás estudiantes, indicando que el estudiante no se motiva por enseñar a los demás los pasos que él siguió para encontrar la solución.

\item En la pregunta N$^{\circ}$ 5 se puede evidenciar que el estudiante opta por esperar que el docente resuelva los problemas y no se motivan a buscar otro camino para resolver el problema, presentando índices de falta de motivación ante algún obstáculo.

\item Se puede evidenciar que en la pregunta N$^{\circ}$ 6 el estudiante considera que la forma de mostrar el contenido de la aplicación se hace de manera creativa, impulsandolo a trabajar sobre el contenido planteado de manera más fluida y simple.

\item El estudiante considera que el apoyo brindado por la aplicación es bastante importante para los procesos de aprendizaje, así, manifestando la necesidad de incluirla en sus estudios.

\end{itemize}


\textbf{Test cognitivo:} Mediante la realización del test por parte de los estudiantes de la Institución Juan María Céspedes Se pudo constatar el perfil cognitivo del estudiante el cual posee características específicas tales como:
\begin{itemize}

    \item El estudiante prefiere visualizar el contenido y los ejercicios de los temas sean de manera gráfica, animada, y didáctica para llevar a cabo sus procesos mentales de enseñanza - aprendizaje según las respuestas a las preguntas N$^{\circ}$ 4, 5, 6, 7, 8, 9, 10, 11, 12, 13, 14 y 15.

    \item El estudiante identifica los elementos de la aplicación de manera separada y conjunta según la respuesta a la pregunta N$^{\circ}$ 2.

    \item El estudiante tiene dificultades para recordar información previamente suministrada de la aplicación, estas dificultades pueden evidenciarse en las preguntas N$^{\circ}$ 1, 3 y 22.

    \item Todos los estudiantes de la muestra fueron clasificados en el perfil dependiente al campo.

    \item El estudiante tuvo dificultades en los cuestionarios que no se visualizaban bien su imagen de respuesta o de pregunta, según las respuestas en la pregunta N$^{\circ}$ 21.
    
    \item El estudiante considera parte importante los recursos de hipermedia visualizados en la aplicación para sus procesos mentales de enseñanza- aprendizaje, según las respuestas en la pregunta N$^{\circ}$ 20.
    
    \item El estudiante considera que el desarrollo de los ejercicios dentro de la aplicación se facilitó con los conocimientos previos y los conocimientos adquiridos en la aplicación, según las respuestas en la pregunta N$^{\circ}$ 19.

    \item El estudiante considera que la estructuración del contenido de los temas y subtemas fueron los adecuados para su perfil y llevar a cabo los procesos mentales de enseñanza- aprendizaje, según las respuestas en la pregunta N$^{\circ}$ 17.

    \item El estudiante se les dificulto los temas que tenían fórmulas complejas e imágenes que no observaban el texto, según las respuestas en la pregunta N$^{\circ}$ 18.

\end{itemize}
